% Paper 5, §1 — Introduction: the layered-defense view
% Pass-C synthesis section. Names the six layers, previews §2-§13, sets up
% the open question ("is the layered-defense framing real or post-hoc?")
% that §14 will answer. Cites only the small canonical anchor set
% (hendrycks2021-mmlu, liang2022-helm, harmbench2024,
% meinke2024-apollo-scheming, greenblatt2024-alignment-faking,
% kirchenbauer2023-watermark, irving2018-debate, sharma2023-sycophancy,
% glazer2024-frontiermath); per-layer primary cites already live in §3-§13.

\section{Introduction: the layered-defense view}
\label{sec:introduction}

A single number does not measure a frontier language model. The
report that ``model $M$ scores $x\%$ on benchmark $B$'' is the
output of a measurement instrument with construct assumptions,
sampling choices, judging procedures, contamination exposure, and
deployment-distribution gaps stacked behind it. None of the
2024--2026 instruments by itself answers the deployment-relevant
question; each one was built because the prior layer left a known
gap; and the gaps left by each new layer are themselves the topic of
the next year's evaluation paper. \textbf{The thesis of this paper
is that the eval-and-alignment landscape is best read as a layered
defense: six measurement layers, each with named failure modes the
next layer is supposed to catch, all of them necessary, none of them
individually sufficient, and the composition between them
imperfect.} Naming the layers and their failure modes is a
precondition for stating, honestly, which deployment-relevant
threats current methods can measure today and which still slip
through.

The six layers, in the order the rest of the paper takes them, are:

\begin{enumerate}
  \item \emph{Knowledge measurement and contamination} --- can the
    model recall what it has been taught, and is the test set
    distinguishable from the training set? Canonical failure mode:
    \emph{measurement-instrument \glsterm{saturation}{saturation} against an answer-\glsterm{key}{key}
    error floor}. The \glsterm{mmlu}{MMLU} lineage \citep{hendrycks2021-mmlu} runs
    into a verifier-error ceiling on the order of single-digit
    percent; reported scores above that ceiling are the
    benchmark's noise floor, not signal
    (\cref{sec:knowledge-benchmarks,sec:contamination}).
  \item \emph{Reasoning-benchmark validity} --- does multi-step
    performance on a fixed exam set transfer to multi-step
    performance off it? Canonical failure mode: \emph{contest-style
    \glsterm{saturation}{saturation}, research-style stall}. AIME and MATH-500 saturate
    while \glsterm{frontiermath}{FrontierMath} \citep{glazer2024-frontiermath} does not, and
    the gap is the gap between in-distribution reasoning-RL gains
    and out-of-distribution generalization
    (\cref{sec:reasoning-benchmarks}).
  \item \emph{Agentic-benchmark realism} --- does end-to-end task
    completion in a fixed environment predict end-to-end task
    completion in a deployed one? Canonical failure mode:
    \emph{realism-versus-measurability tradeoff}. SWE-Bench, OSWorld,
    GAIA, and $\tau$-Bench occupy four points on a single curve
    where greater realism buys reliability cost (versions drift,
    websites change, judges disagree); reported scores across years
    are not directly comparable without environment-version lock
    (\cref{sec:agentic-benchmarks}).
  \item \emph{Adversarial robustness} --- under a search-based
    attacker, does the deployed system keep refusing the behaviors
    its training said it should refuse? Canonical failure mode:
    \emph{\glsterm{asr}{ASR} non-comparability}. The four canonical
    attack-success-rate (\glsterm{asr}{ASR}) instruments
    \citep{harmbench2024} use different target behaviors,
    attacker LLMs, judges, and success criteria; the same model's
    $5\%$ \glsterm{asr}{ASR} on one and $30\%$ \glsterm{asr}{ASR} on another is not a
    contradiction, it is two different constructs
    (\cref{sec:adversarial-robustness}).
  \item \emph{Alignment-threat detection} --- does the model behave
    aligned because it \emph{is} aligned, or because it has
    capability and incentive to look aligned during evaluation?
    Three sub-layers compose here: \glsterm{sycophancy}{sycophancy}
    \citep{sharma2023-sycophancy} as a measured \glsterm{rlhf}{RLHF} artifact
    (\cref{sec:sycophancy}); \glsterm{in-context-scheming}{in-context scheming}
    \citep{meinke2024-apollo-scheming} as an elicited capability
    finding (\cref{sec:scheming}); alignment-faking
    \citep{greenblatt2024-alignment-faking} as strategic compliance
    under perceived training-vs-deployment shift
    (\cref{sec:alignment-faking}). Canonical failure mode for the
    layer as a whole: \emph{capability under elicitation does not
    imply propensity in deployment}, and the field has no
    deployment audit that closes the gap.
  \item \emph{Provenance and \glsterm{scalable-oversight}{scalable oversight}} --- once outputs
    leave the model, can their origin be certified, and can a
    bounded human resource budget align a model whose capability
    already exceeds the overseer's? Watermarking
    \citep{kirchenbauer2023-watermark} addresses the first
    (\cref{sec:watermarking}); \glsterm{debate}{debate}
    \citep{irving2018-debate} and its sibling protocols address
    the second (\cref{sec:scalable-oversight}). Canonical failure
    mode: \emph{robustness-under-removal-pressure}. The
    cryptographic part of watermarking is solved; surviving
    paraphrase, fine-tuning, and collusion is not. The theoretical
    part of \glsterm{scalable-oversight}{scalable oversight} is solved (\glsterm{debate}{debate} decides
    $\mathbf{PSPACE}$ under polynomial-time judging); demonstrating
    bounded-judge honesty on open-ended domains is not.
\end{enumerate}

\begin{intuition}
The picture is a stack, not a benchmark. Each layer is an
instrument the field built because the prior layer left a known
hole: contamination detection because knowledge benchmarks could be
gamed by leakage; reasoning benchmarks because knowledge benchmarks
saturated; agentic benchmarks because reasoning benchmarks stayed
text-only; adversarial robustness because agentic capability without
refusal training is dangerous; alignment-threat evals because
refusal can mask misaligned cognition; provenance and oversight
because measurement of \emph{current} models does not yet
generalize to \emph{superhuman} ones. Returning to canonical form:
each layer pairs (a) a construct it claims to measure with (b) a
named failure mode that motivates the next layer. The
layered-defense framing is itself an intuition, and \cref{sec:contradictions} grounds it
by enumerating, layer by layer, the actual instruments the field
has shipped --- so the framing is checkable, not just rhetorical.
\end{intuition}

\subsection*{The structural commitment}

This paper commits to three claims about each layer
simultaneously, and the body sections discharge those commitments
one layer at a time:

\begin{itemize}
  \item \emph{Necessary.} Drop any one layer and a real
    deployment-relevant threat goes unmeasured. Without
    contamination detection, knowledge benchmarks read leakage
    as capability. Without reasoning-benchmark validity, \glsterm{rlvr-3}{RLVR}-trained
    models look generally smarter than they are. Without agentic
    benchmarks, end-to-end deployment risk is invisible. Without
    adversarial robustness, refusal training is unfalsified.
    Without alignment-threat detection, a model that schemes well
    enough to pass scores cleanly. Without provenance and
    oversight, even a perfectly-aligned-on-current-distribution
    model leaves the deployment-distribution shift unaddressed.
  \item \emph{Insufficient.} No single layer answers the
    deployment-relevant question on its own. A high \glsterm{mmlu}{MMLU} score
    that survives \glsterm{mmlu-redux-2}{MMLU-Redux} still says nothing about agentic task
    completion or \glsterm{scheming}{scheming} under elicitation. A low \glsterm{harmbench-2}{HarmBench} \glsterm{asr}{ASR}
    says nothing about whether the model has dangerous capability
    in the first place. A clean Apollo-\glsterm{scheming}{scheming} evaluation says
    nothing about non-elicited propensity. The instruments measure
    \emph{what they measure}, and the construct-validity question
    --- whether what they measure is what we want to know --- is
    open at every layer.
  \item \emph{Imperfectly composed.} The layers do not commute.
    A model that scores well on $\S$3 (knowledge) but fails $\S$11
    (alignment-faking) is a different deployment risk than one
    that fails $\S$3 and passes $\S$11. The field reports each
    layer's score in isolation; the joint distribution of scores
    across the stack is what determines deployment risk, and that
    joint distribution is reported in essentially no published
    eval.
\end{itemize}

\subsection*{What this paper is and is not}

\paragraph{What it is.} A monograph on the 2024--2026 state of \glsterm{llm}{LLM}
evaluation and alignment-threat detection, organized by measurement
layer rather than by benchmark or by lab. Every load-bearing
empirical claim --- benchmark score, \glsterm{asr}{ASR}, contamination percentage,
\glsterm{sycophancy}{sycophancy} rate, \glsterm{scheming}{scheming}-suite hit-rate, alignment-faking rate,
\glsterm{watermark}{watermark} detection threshold --- cites a primary source whose
excerpt the KB has tracked down to a section anchor. Every formal
construction --- the IRT 2PL form (\cref{sec:measuring-framing}),
\glsterm{helm}{HELM}'s scenario-by-metric matrix (\cref{sec:measuring-framing}),
\glsterm{asr}{ASR}'s set-cardinality definition (\cref{sec:adversarial-robustness}),
the green-list partition function (\cref{sec:watermarking}), the
\glsterm{debate}{debate} game's $\mathbf{PSPACE}$ construction
(\cref{sec:scalable-oversight}) --- is transcribed verbatim from the
cited paper, in the notation that paper uses. Every cross-source
disagreement is tagged \texttt{[CONTRADICTION]} in the body and
gathered for direct discussion in \cref{sec:contradictions}.

\paragraph{What it is not.} A complete benchmark catalogue, a
mitigation manual, or a regulatory roadmap. The paper is
empirical-claim-dense rather than mechanism-dense; the field is
more empirical than the ones Papers~1--3 cover, and the math layer
is correspondingly thinner. Where the load-bearing 2024--2026
demonstration is by a single research group with no independent
replication --- the alignment-faking setup
\citep{greenblatt2024-alignment-faking}, the \glsterm{scheming}{scheming} suite
\citep{meinke2024-apollo-scheming}, \glsterm{frontiermath}{FrontierMath}
\citep{glazer2024-frontiermath} --- we say so explicitly rather
than papering over the single-source status. Mitigations and
deployment-policy proposals are referenced where they exist
(Responsible Scaling Policies, constitutional methods, scalable
oversight protocols) but are not the paper's subject. Earlier
papers in this series carry the prerequisite machinery: Paper~2's
pre-training-data and reasoning-RL stages condition the inputs to
\cref{sec:knowledge-benchmarks,sec:contamination,sec:reasoning-benchmarks};
Paper~4's probing and \glsterm{circuit}{circuit}-tracing methods are a complementary
attack surface for
\cref{sec:scheming,sec:alignment-faking}; Paper~3's
process-supervision and Paper~2's weak-to-strong generalization
results bound \cref{sec:scalable-oversight}. We name those threads
where they enter and do not re-derive them.

\subsection*{The open question this paper does and does not answer}

A natural objection --- which we take seriously --- is that the
six-layer framing might be \emph{post-hoc taxonomy}: a
narratively-clean partition imposed on a messy empirical
literature, where the actual instruments do not in fact correspond
to independent layers and where ``layered defense'' is a metaphor
that pre-commits to more structure than the field has built. We
address this question directly in \cref{sec:contradictions} by
enumerating, layer by layer, the actual published instruments
(\glsterm{mmlu}{MMLU} and \glsterm{mmlu-redux-2}{MMLU-Redux} for layer 1; the contamination survey
detection methods for the contamination sub-layer; AIME, MATH-500,
GPQA, \glsterm{frontiermath}{FrontierMath}, \glsterm{hle}{HLE} for layer 2; SWE-Bench, OSWorld, GAIA,
$\tau$-Bench for layer 3; \glsterm{harmbench-2}{HarmBench}, \glsterm{jailbreakbench-2}{JailbreakBench}, Crescendo,
Wei~2023 for layer 4; the four \glsterm{sycophancy}{sycophancy} datasets, Apollo's six
\glsterm{scheming}{scheming} evaluations, Greenblatt's free-tier/paid-tier setup for
layer 5; Kirchenbauer's green-list construction, \glsterm{debate}{debate}, \glsterm{iterated-amplification}{IDA},
recursive reward modeling, weak-to-strong for layer 6) and showing
that each has at least one \emph{independently developed} measurement
artifact and at least one \emph{independently named} failure mode.
The framing is real if and only if both conditions hold; we do not
preview the answer here, only flag that the question is the
load-bearing one for the paper's structural claim.

\subsection*{Roadmap}

\Cref{sec:measuring-framing} fixes the eval-methodology \glsterm{vocabulary}{vocabulary}
the rest of the paper inherits: validity, reliability, construct,
contamination, ecological validity; the three orthogonal axes
(capability vs alignment; knowledge vs reasoning vs agency; static
benchmark vs adversarial \glsterm{probe-2}{probe} vs deployment audit); \glsterm{helm}{HELM}'s
scenario-by-metric matrix \citep{liang2022-helm} as the closest the
field has to a theory of measurement; the 2PL IRT form as a
calibration tool. \Cref{sec:knowledge-benchmarks,sec:contamination}
walk the knowledge layer and its contamination boundary.
\Cref{sec:reasoning-benchmarks} takes the reasoning-validity layer
and locates the contest-vs-research split.
\Cref{sec:agentic-benchmarks} takes the agentic layer and the
realism-vs-measurability tradeoff.
\Cref{sec:adversarial-robustness,sec:dangerous-capability} cover
the two safety-evaluation sub-layers (refusal robustness and
underlying capability).
\Cref{sec:sycophancy,sec:scheming,sec:alignment-faking} cover the
three alignment-threat sub-layers, ordered by increasing
strategic-cognition requirement on the model.
\Cref{sec:watermarking} takes the provenance instrument.
\Cref{sec:scalable-oversight} takes the oversight question.
\Cref{sec:contradictions} concentrates the
\texttt{[CONTRADICTION]} markers from the body sections, names
what evidence would close each one, and returns to the
post-hoc-framing question. The closing claim is that the layered
defense holds today, but each layer's instruments are aging at
different rates; the next round of evaluation infrastructure has
to attack the specific gaps we name, not produce more of what we
already have. We turn first to the framing:
\cref{sec:measuring-framing}.
